\Section{Encoding}
\label{sec:enc}

We describe the implementation of function values as a translation into a conceptual (pseudo-code style) imperative intermediate verification language (IVL) inspired by Boogie~\cite{boogieIVL,thisisboogie2}.
The encoding is given at the IVL level because this is where the interesting design decisions become visible --- the subsequent step to pure SMT is straightforward and uses standard techniques.
The fragment of the IVL used in this section is familiar: procedures with $\HAVOC$, $\ASSUME$, $\ASSERT$, pure functions, \emph{algebraic datatypes} (tagged unions with selectors, like Rust enums), universal \emph{axioms} with \emph{triggers}, and polymorphic map types.
State is represented by Boogie-style maps; in particular, each resource type $\tau$ yields a map $M_\tau \in \Addr \rightarrow \tau$, with $\Addr$ being the domain of addresses in Move (256-bit integers).
The representation is similar to that in the earlier \MVP paper~\cite{DillGPQXZ22}.

In the following we focus on the aspects specific to function values: the representation of function values (Sec.~\ref{sec:enc-fvals}), the memory access and modification sets used to define the invocation schema (Sec.~\ref{sec:enc-mem}), the encoding of invocation of function values (Sec.~\ref{sec:enc-invoke}), the encoding of behavioral predicates via axioms (Sec.~\ref{sec:enc-bp}), and finally the encoding of intermediate state labels (Sec.~\ref{sec:enc-frame}).
Everything in this section operates \emph{after} monomorphization, so type parameters have been specialized and the set of closure, parameter, and field variants reaching a given function type is known statically~\cite{DillGPQXZ22}.

\SubSection{Representing Function Values}
\label{sec:enc-fvals}
For each function type $\tau = (T_1,\ldots,T_n) \to T$ that occurs in the monomorphized program, we create a datatype in the IVL which represents all the potential \emph{sources} of function values in a given program.
To this end, \MVP's mono-analysis records three sets: the set $\Clo_\tau$ of concrete functions $f$ constructing closures of type $\tau$, the set $\Par_\tau$ of function-typed parameters of verification targets, and the set $\Fld_\tau$ of storable function-typed struct fields whose type is $\tau$.
The union of these three sets is materialized as an algebraic datatype in the IVL --- one constructor per variant:

\begin{ivl}
\DATATYPE \tau\ \{\\
\t1 C_{f}(\sq{c : T}), \ldots, P_{f,x}(), \ldots, F_{\sigma.x}(n : \nat), \ldots \\
\}
\end{ivl}

\noindent Here $\sq{c:T}$ are the types of the values captured by closure $f$; $P_{f,x}$ is the parameter variant for function-value parameter $x$ of function $f$; and $F_{\sigma.x}$ is the field variant for struct field $x$ of struct type $\sigma$ (we use $\sigma$ here to avoid confusion with the state-label variable $S$ introduced in Sec.~\ref{sec:enc-frame}).
Parameter variants are nullary because, in a well-typed verification condition (VC), the value of an entrypoint function-typed parameter is a skolem constant pinned to its variant at entry; no runtime payload is required.
Field variants carry an integer discriminator $n$ so that two locations in memory holding distinct field values are not forced to be equal --- a datatype with a single nullary constructor would otherwise collapse all field values to one.

Notice that this representation captures \emph{all} possible functions of type $\tau$ relevant for a given VC.
So we do not actually need to deal with an arbitrary number of function values when it comes to higher-order function verification, but only those which can be derived from the program in a given context.
The representation is further narrowed if the function value of a given type is not stored in a structure which ends in global memory.
Storing a function value in a struct field introduces dynamic dispatch: two different field locations may hold different closures, and the discriminator $n$ in $F_{\sigma.x}(n)$ ensures the solver treats them as distinct.
Nevertheless, we are able to pinpoint a specification to a particular $F_{\sigma.x}(n)$, as illustrated by the AMM example (Sec.~\ref{sec:amm}).


\SubSection{Computing Memory Accesses}
\label{sec:enc-mem}

Invoking a function value requires reasoning about which memory maps a given function type $\tau$ may read ($\Mem_\tau$) and which locations it may modify ($\Mod_\tau(\sq{p})$); we define both here before presenting the invocation schema.
Intuitively, the access set is the \emph{read footprint} of the function values of type $\tau$ --- it delimits which parts of memory their behavior may depend on --- while the modification set is their \emph{write footprint}, providing the frame condition: everything outside of it is unchanged by an invocation.
To recall, in the Move language, persistent global storage, or memory, can be accessed via type indexing.
For instance, if !R! is a resource type, we write !&mut R[addr]! in the language to obtain a mutable reference to the memory which can then be subsequently mutated.

\MVP's representation of memory uses type-indexed maps $M_\rho \in \Addr \fun \rho$, with $\rho$ a resource type (a !struct! in Move with the !key! ability).
In the imperative procedures of the IVL, those maps are implicitly available as global variables.
But in pure expressions like behavioral predicates, they need to be made explicit as parameters; this is why $\Mem_\tau$ appears as an explicit argument to behavioral predicates throughout this section.

A few definitions for memory maps are needed.
Note that these are definitions on the meta level of the description here and not part of the IVL; they are used to describe pseudo code in terms of the IVL.

First, with $\Mem_R = \{ M \in \Addr \fun \rho \mid \rho \in R\}$ we denote a set of memory maps for the resource types in $R$.
The $R$ can be omitted if it is clear from the context. For convenience, we write $\Type(\Mem)$ for the set of all types $\rho$ used in the memory maps of $\Mem$.
With $\Mem_R \downarrow S$, where $S \subseteq R$, we denote the projection $\{ M \in \Mem_R \mid \rho(M) \in S \}$.

In addition to memory, we also need to reason about modifications.
A modification is a set of pairs of a type and an expression of the IVL of type $\Addr$, $\Mod \in \power (\Type \cross \Expr)$.
We write $\Mod(\sq{p})$ for a set of modifications built on expressions referencing parameters in $\sq{p}$.
For instance, the declaration |modifies_of<f>(a: address) R[a]| (introduced below) contributes the pair $(R, a)$; at an invocation of !f!, $a$ is bound to the actual argument, so $\Mod(\sq{p})$ gives the modified locations as a function of the arguments $\sq{p}$.


\Paragraph{General Access}
Memory access $\Mem_\tau$ is computed as the union of all the accesses in the variants of $\tau$, $\Mem_{C_f} \cup \Mem_{P_{f,x}} \cup \Mem_{F_{\sigma.x}}$.
For a closure, $\Mem_{C_f}$  can be statically derived from the code.
For parameters and fields, the information is declared in the spec block:

% @formatter:off
\begin{MoveBox}
  spec foo(f: |T|S, g: |T|S) {
    reads_of<f> R; // f can read resource R
    reads_of<g> *; // g can read any resource
                   // in scope of the VC
  }
\end{MoveBox}
% @formatter:on

\noindent For the wildcard operator !*!, all memory which is read or written as part of the currently generated VC is included, plus an abstract memory map for any `unknown' memory the function may depend on.
The unknown memory is needed for a sound encoding.
For instance, if we verify function !foo! above in a context where no memory is involved anywhere, we must not assume that !g(x)! delivers the same value in different states, as it may depend on this unknown additional memory.
However, we can still know that the specification of !g! holds in every possible state.

\Paragraph{Write Accesses}
Similarly to general access, the modifies set $\Mod_\tau$ is computed as the union of all the modifications in the variants of $\tau$, $\Mod_{C_f} \cup \Mod_{P_{f,x}} \cup \Mod_{F_{\sigma.x}}$.

For closures, $\Mod_{C_f}$ is derived from the code and, if present, the !modifies R[a]! declaration in the language.
These existed before higher-order functions, for specifying what functions modify.

For parameters and fields, the information is declared in the spec block, where for the !modifies_of<f>! clause, the parameters of the function value itself are passed:

% @formatter:off
\begin{MoveBox}
  spec foo(f: |address|) {
    // f can modify resource R[a]
    modifies_of<f>(a: address) R[a];
  }
\end{MoveBox}
% @formatter:on

\noindent If no !modifies_of! is declared, it is assumed that the function does not modify global memory.

Notice that a wildcard (!modifies_of<f> *!) is not supported, since the invocation of !f! would render the memory arbitrary (no effective frame condition).
Consider a generic utility !for_each_account(f: |address|)! applying a caller-supplied, state-mutating callback: one would like to specify it once, for callbacks with open-ended write footprints, but today the footprint must be declared per resource, as in !modifies_of<f>(a) R[a]!.
Lifting this restriction is left for future work.

\SubSection{Invoking Function Values}
\label{sec:enc-invoke}

When a function value is invoked, we are looking at either a concrete closure value $C_{g}$ or a parameter or field selection as described above.
In the former case, one can simply delegate to calling the actual underlying function $f$, by composing the captured and provided arguments into a full argument list.

In the other cases of parameters and field selections, we use the behavioral predicates as defined for the function value, using the standard schema for encoding a call to a Move function via the function's specification.
This reduces the problem of invocation to how those behavioral predicates are encoded.
In the procedure below, $\Mem_\tau$ denotes the projection of the current global memory state to the resource types in $\tau$'s access set (as computed in Sec.~\ref{sec:enc-mem}).
\begin{ivl}
\PROC invoke_\tau(x: \tau, \sq{p:T}) \RETURNS (\sq{r: T}) \\
\t1 \IF x \IS\ C_{f}(\sq{c}) \\
\t2    \sq{r} := \CALL f(\sq{c}, \sq{p}) \\
\t1 \ELSE \\
\t2    \ASSERT requires_\tau(x, \Mem_\tau, \sq{p}) \\
\t2    \IF\ aborts_\tau(x, \Mem_\tau, \sq{p}) \\
\t3       abort\_flag := \TRUE \\
\t2    \ELSE \\
\t3       \LET \Mem'_\tau := \Mem_\tau \\
\t3       \HAVOC \Mod_\tau(\sq{p}),\ \bar{r} \\
\t3       \ASSUME ensures_\tau(x, \Mem'_\tau, \Mem_\tau, \sq{p}, \sq{r}) \\
\end{ivl}

\noindent Here $\LET \Mem'_\tau := \Mem_\tau$ freezes the pre-invocation memory into $\Mem'_\tau$, while $\Mem_\tau$ becomes the post-invocation state after $\HAVOC$.
We therefore pass $\Mem'_\tau$ as the pre-state and $\Mem_\tau$ as the post-state to $ensures_\tau$ (note that $\Mem'$ consistently denotes the pre-state throughout this section).
$\HAVOC \Mod_\tau(\sq{p})$ is shorthand for havocing, for each $(\rho, a) \in \Mod_\tau(\sq{p})$, the entry $M_\rho[a]$ of the type-indexed memory map for $\rho$.
$abort\_flag$ is a prover-internal boolean global that signals an abort to the enclosing verification condition, consistent with how \MVP models Move's abort semantics throughout the IVL.
The $\IF$ branch reduces to a standard function call $\CALL f(\sq{c}, \sq{p})$, which \MVP encodes using the same $\ASSERT$/$\HAVOC$/$\ASSUME$ schema as the $\ELSE$ branch, specialized to $f$'s declared spec.
When $f$ is opaque, the two branches are in fact semantically equivalent: an opaque call is already compiled to this schema driven by $f$'s spec, matching what the $\ELSE$ branch does for the abstract variant.
The $\ELSE$ branch therefore introduces no new encoding machinery; it applies the existing opaque-call schema, driven by the behavioral predicates of the unknown variant $x$ rather than a concrete function's spec.

\SubSection{Behavioral Predicates}
\label{sec:enc-bp}

Behavioral predicates are derived by switching over the variants of the function value representation $\tau$, similarly to the $invoke_\tau$ procedure above, delegating extraction of the underlying predicates to per-variant functions.
Intuitively, each behavioral predicate is total over the variants of $\tau$: applied to a closure, it unfolds to the underlying function's specification; applied to an abstract parameter or field variant, it remains an opaque symbol of which only the user-declared conditions are known.

Below, the definition for ensures is given.
The other predicates are defined similarly.
Specific to ensures is that it needs to take care of the frame condition: the caller havoced the full modification set, but only a subset of that is constrained by the given variant.
The handling needs to account for the fact that modification addresses are symbolic expressions that may alias.

In what follows, we write $\mathit{pred}_\tau[\Mem', \Mem](\ldots)$ with square brackets for the memory pair: $\Mem'$ is the pre-state and $\Mem$ the post-state, following the convention established in Sec.~\ref{sec:enc-invoke}.
These memory arguments must be passed explicitly to pure functions, whereas they are implicit globals in IVL procedures.
We write $\Mod_x$ for the modification set of the specific variant of $x$ --- that is, $\Mod_{C_f}$ when $x = C_f(\sq{c})$, $\Mod_{P_{g,y}}$ when $x = P_{g,y}()$, and $\Mod_{F_{\sigma.x}}$ when $x = F_{\sigma.x}(n)$.
We write $\Mem'(\rho)$ for the map $M_\rho$ in the pre-state, and $\Mem(\rho)$ for the map $M_\rho$ in the post-state.

\begin{ivl}
  \FUN ensures_\tau[\Mem', \Mem](x: \tau, \sq{p: T}, \sq{r: U}): \BOOL \\
\t1  \bigwedge_{(\rho, a) \in \Mod_\tau}( \\
\t2    \IF \rho \notin \Type(\Mod_x) \\
\t3       \Mem(\rho) = \Mem'(\rho) \\
\t2    \ELSE \\
\t3       (\bigwedge_{(\rho, a') \in \Mod_x} a' \neq a) \\
\t4          \implies \Mem(\rho)[a] = \Mem'(\rho)[a] \\
\t1  ) \land ensures_x[\Mem'_x, \Mem_x](p, r)
\end {ivl}

The relational representation of a function via ensures does not yet capture that Move functions are \emph{deterministic} (Sec.~\ref{sec:move}): for a fixed function value, arguments, and pre-state, there is exactly one outcome.
Functional behavior is established by introducing $result_\tau$ as an uninterpreted function, connected to ensures via an axiom (from here on, quantifiers are annotated with their \emph{trigger} in braces, further discussed in Sec.~\ref{sec:imp-notes}):
\begin{ivl}
  \FUN result_\tau[\Mem', \Mem](x: \tau, \sq{p: T}) \fun \sq{U} \\
  \AXIOM \forall x, \sq{p} @ \{result_\tau[\Mem', \Mem](x, \sq{p})\} \\
\t1 \LET \sq{r} := result_\tau[\Mem', \Mem](x, \sq{p}) @ \\
\t2 ensures_\tau[\Mem', \Mem](x, \sq{p}, \sq{r})
\end {ivl}

The result tuple $\sq{U}$ comprises the declared results of $\tau$ and, per mutable-reference parameter, the post-state of the referenced value; thus both the explicit results and the implicitly modified local state are pinned as a function of the function value, its arguments, and memory.
For a resource which the variants of $\tau$ only read, a single memory map is passed, so the outcome is determined by the pre-state alone, as determinism demands; the post-state of \emph{modified} global memory remains relationally constrained by the invocation schema (havoc within the frame $\Mod_\tau$, then assuming ensures).

The axiom is deliberately one-directional, stating that the tuple chosen by $result_\tau$ satisfies $ensures_\tau$.
The seemingly natural biconditional $ensures_\tau(x, \sq{p}, \sq{r}) \iff \sq{r} = result_\tau(x, \sq{p})$ would force \emph{all} solutions of ensures to coincide --- inconsistent as soon as the declared ensures of some variant under-constrains its outcome, such as a mere |ensures result > 0|.

The actual conditions extracted depend on the variant.
For closures, they are directly extracted from the spec block.
For parameters and fields, we introduce uninterpreted functions; for ensures specifically, we introduce an uninterpreted per-variant result function and \emph{define} the variant's ensures as its graph, so functional behavior holds by construction (the field variant is analogous to the parameter variant shown, with the discriminator $n$ as additional argument):

\begin{ivl}
  \FUN ensures_{C_f(\sq{c})}[\Mem', \Mem](\sq{p: T}, \sq{r: U}): \BOOL \\
\t1    \langle \textit{extract from spec block of}~~f(\sq{c}, \sq{p}) \rangle \\
  \FUN result_{P_{f,x}}[\Mem', \Mem](\sq{p: T}) \fun \sq{U} \\
  \FUN ensures_{P_{f,x}}[\Mem', \Mem](\sq{p: T}, \sq{r: U}): \BOOL \\
\t1    \sq{r} = result_{P_{f,x}}[\Mem', \Mem](\sq{p}) \\
\end{ivl}

\noindent The meaning of the uninterpreted functions is determined by how they are injected into a VC via the regular mechanism by which pre-/postconditions and invariants are handled.
Consider a precondition like |requires !aborts_of<param>(x)|, where !param! is a function-value parameter of a function !foo!.
This will be injected in the VC as $\ASSUME \lnot aborts_{P_{foo, param}}(x)$.
Similarly, for a data invariant of a field, an assumption will be generated when the field is selected.

\Paragraph{Compliance checking}
The same mechanism yields the compliance obligations of Sec.~\ref{sec:amm} when a concrete function value flows into an abstract position.
Since behavioral predicates are first-order, compliance verification mimics the conventional handling of pre-/postconditions: the condition |requires !aborts_of<param>(x)| above, assumed on the definition side of !foo!, is \emph{asserted} at call sites of !foo!, where !param! is bound to an actual function value --- say a closure $C_g(\sq{c})$ --- so the assertion reduces, via the per-variant extraction, to a condition over $g$'s specification.
Data invariants over function-valued fields, as in the Pool example, create the same obligations at pack time; constraints via $result_\tau$, including relational ones, are discharged through the result axiom.
The resulting VCs quantify over the function value's arguments (and possibly memories) and may involve non-linear arithmetic (Sec.~\ref{sec:amm}); they are inherently harder than application-side VCs and may need proof hints.


\SubSection{Intermediate State Labels}
\label{sec:enc-frame}

Recall from Sec.~\ref{sec:move-hof} that a state label !S! denotes a program point.
A specification mentioning !S! talks about a state \emph{in the middle} of a function's execution, invisible in a pre-/post-state contract; the encoding materializes it as an existentially quantified snapshot, connected to both ends by the conditions mentioning !S!.
Formally, a \emph{state} at label $S$ is a pair $(\Mem^S, \bar{v}^S)$ consisting of a memory snapshot and a tuple of local values.
The memory component $\Mem^S$ has the same type as the ambient memory: a tuple of type-indexed maps $M_\rho \in \Addr \to \rho$ for each resource type $\rho$ in scope at that program point, quantified existentially in the IVL.
The local component $\bar{v}^S$ captures values of mutable-reference parameters at the intermediate point: for a closure over a mutable reference of type $T$ with no global memory effects, $\bar{v}^S$ is a single value $x^S : T$, and $\Mem^S$ is vacuous (the function touches no global resources).
When a closure both accesses global memory and takes mutable references, both components are non-trivial.
Here we focus on the global-memory case (vacuous $\bar{v}^S$); Sec.~\ref{sec:imp-notes} notes how the mutable-reference case is handled.

Consider the following specification using intermediate state labels:

% @formatter:off
\begin{MoveBox}
  struct R(u64) has key;
  spec foo(f: |address|, g: |address|, a: address) {
    ensures ..S |~ ensures_of<f>(a)
    ensures S |~ R[a].0 > 0;
    aborts_if S |~ aborts_of<g>(a);
    ensures S.. |~ ensures_of<g>(a)
  }
\end{MoveBox}
% @formatter:on

As outlined in the last section, we want to be able to extract behavioral predicate functions from this specification.
For the ensures we can come up with the following, which simply replicates part of the spec:

\begin{ivl}
  \FUN ensures_{C_{foo}}[\Mem', \Mem](a) \\
\t1  \exists \Mem^S @
      \<ensures_\tau(f,\Mem',\Mem^S,a) \land R[\Mem^S][a].0 > 0 \land \\
        ensures_{\tau'}(g,\Mem^S,\Mem,a)\>
\end{ivl}

For the aborts which happens at the intermediate state !S!, the state !S! needs to be first established before the aborts condition can be checked for !g!.
To achieve this, all ensures conditions which reach !S! are combined with the aborts condition, as in:

\begin{ivl}
  \FUN aborts_{C_{foo}}[\Mem', \Mem](a) \\
\t1  \exists \Mem^S @
      \<ensures_\tau(f,\Mem',\Mem^S,a) \land R[\Mem^S][a].0 > 0 \land \\
        aborts_{\tau'}(g,\Mem^S,a)\>
\end{ivl}

This is semantically sound because we can assume inductively that the function $foo$ has successfully verified. For a successfully verified function, where $\sq{A}$ are the |aborts_if| conditions and $\sq{P}$ the |ensures| conditions, it holds:
$$
   \bigvee \sq{A} \lor \bigwedge \sq{P}
$$

\noindent That is, either the function aborts under one of the given conditions, or all of its ensures conditions hold.

In general, since the relation between state labels spawned by predicates of the form $\slabeled{S_1..S_2}{p}$ must be acyclic, one can extract from $\sq{P}$ those conditions which lead into $\Mem^S$, and from $\sq{A}$ those which start from state $\Mem^S$.
Given this, the aborts condition can be synthesized using the following scheme:
$$
    \bigwedge \sq{P}_{..\Mem^S} \land (\bigvee \sq{A}_{\Mem^S} \lor
    \bigwedge \sq{P}_{\Mem^S..\Mem^T} \land (\bigvee \sq{A}_{\Mem^T} \lor
    \bigwedge \ldots ))
$$

\noindent Here, $\sq{P}_{..\Mem^S}$ establish state $\Mem^S$ in which we check for the aborts condition $\sq{A}_{\Mem^S}$,
or recurse over the remaining $\sq{P}$ and $\sq{A}$.

\SubSection{Implementation Notes}
\label{sec:imp-notes}

A straightforward implementation following the conceptual description above leads to SMT timeouts on non-trivial examples.
The root cause is that the presentation writes behavioral predicates as IVL functions with explicit bodies ($\FUN \ldots = \mathit{body}$), and Boogie unfolds function bodies eagerly: any quantifier inside a body --- common in arithmetic identities or prelude vector axioms --- spawns trigger instantiations at every call site, causing combinatorial blowup in the SMT search.

We overcame this by adopting a pattern that decouples the \emph{definition} of behavioral predicates from their \emph{use}: declare each predicate as an \emph{uninterpreted} function and connect it to the per-function specification via axioms with explicit, carefully chosen triggers.
This pattern is not specific to higher-order functions; it is a general Boogie encoding discipline for modular specifications that keeps SMT instantiation under control.
The remaining implementation choices are natural refinements of applying this pattern consistently:

\begin{Itemize}
    \item \emph{Behavioral predicates are uninterpreted functions connected by axioms, not Boogie functions with explicit bodies.}
    The implementation declares $ensures_\tau$ as an uninterpreted Boogie function and connects it to the per-function spec via an axiom of the form $(\forall\, \ldots\; ::\;  \mathit{eval\_call} \iff \mathit{rhs})$ with an explicit trigger pattern.
    Axioms with explicit triggers fire only when the trigger pattern matches a ground term in the proof context, keeping instantiation under control.

    \item \emph{Per-variant trigger specialization.}
    The three variant kinds receive structurally different axioms.
    Shown for $ensures_\tau$ with the memory parameters elided for brevity, the axiom schemas are:

% The zed environment derives its width from \hsize, which lists do not
% reduce; the minipage makes it respect the Itemize indentation.
\begin{minipage}{\linewidth}
\begin{ivl}
\AXIOM \forall \sq{c}, \sq{p}, \sq{r} @ \{ensures_\tau(C_f(\sq{c}), \sq{p}, \sq{r})\} \\
\t1 ensures_\tau(C_f(\sq{c}), \sq{p}, \sq{r}) \iff ensures_{C_f(\sq{c})}(\sq{p}, \sq{r}) \\
\also
\AXIOM \forall \sq{p}, \sq{r} @ \{ensures_\tau(P_{f,x}(), \sq{p}, \sq{r})\} \\
\t1 ensures_\tau(P_{f,x}(), \sq{p}, \sq{r}) \iff ensures_{P_{f,x}}(\sq{p}, \sq{r}) \\
\also
\AXIOM \forall f, \sq{p}, \sq{r} @ \{ensures_\tau(f, \sq{p}, \sq{r})\} \\
\t1 f \IS F_{\sigma.x} \implies (ensures_\tau(f, \sq{p}, \sq{r}) \\
\t2 \iff ensures_{F_{\sigma.x}}(f.n, \sq{p}, \sq{r}))
\end{ivl}
\end{minipage}

    For closure variants, the trigger is the evaluator application \emph{containing} the constructor term $C_f(\sq{c})$, with the captures bound by the quantifier --- a bare constructor pattern would not cover the remaining bound variables; for parameter variants, the nullary ground constructor $P_{f,x}()$ plays the same role.
    For struct-field variants no such constructor term exists: at call sites a field-stored function value is an opaque datatype term loaded from memory.
    The trigger is therefore the evaluator application itself, with the variant pinned by the $f \IS F_{\sigma.x}$ guard and the discriminator passed via the selector $f.n$.
    In each case the axiom is instantiated only where the variant --- or, for fields, the evaluator application --- is syntactically present in the proof context.

    \item \emph{Constraining state-label conjuncts are dropped from the existential.}
    The conceptual behavioral predicate for an !aborts_if! at an intermediate state $\Mem^S$ includes \emph{all} ensures-with-$S$ conjuncts that reach $\Mem^S$ (including constraining ones such as $R[\Mem^S][a].0 > 0$).
    The implementation restricts the existential body to the \emph{defining} conjuncts --- those introduced by label-defining operations like $\ensuresof{f}{\bar{x}}$, $\resultof{f}{\bar{x}}$, or the publish/remove/update builtins.
    By the validity invariant $\bigvee \sq{A} \lor \bigwedge \sq{P}$, the constraining conjuncts hold at the witness chosen during verification of !foo!, so dropping them is sound and keeps the existential body small.

    \item \emph{State-label existential is flat, not recursive.}
    The recursive scheme above is logically equivalent to a single flat $(\exists\, \Mem^S, \Mem^T, \ldots\; ::\;  \ldots)$ wrapping the conjunction of defining fragments and the kind-specific clauses.
    The implementation uses the flat form, reducing the size of emitted Boogie and avoiding the additional SMT quantifier nesting that the recursive scheme would introduce.
\end{Itemize}
